\chapter{BIBLIOGRAFIE. REGULI DE CITARE}

\section{Bibliografie}

Lucrările se scriu în cadrul referințelor bibliografice în ordinea alfabetică a numelor primilor autori, numerotându-se cu cifre arabe între paranteze drepte. Titlul lucrării se scrie în referința bibliografică exact cum este tipărit în publicația citată.

Mai jos sunt prezentate câteva exemple de scriere a referințelor bibliografice – primul exemplu pentru cărți \cite{Burileanu2022}, al doilea pentru capitole de cărți \cite{Kao2020}, al treilea pentru lucrări de conferință \cite{Zhang2021}, al patrulea pentru articole (în jurnale/reviste) \cite{Yi2022}, iar ultimele două pentru resurse disponibile online \cite{Zeiler2012,MATLAB2025}.

Trimiterile din text la bibliografie se fac scriind drept, între paranteze drepte, numărul de ordine al lucrării/lucrărilor – ex.: \cite{Burileanu2022}; \cite{Kao2020,Zeiler2012}; \cite{Zhang2021,Yi2022,Zeiler2012,MATLAB2025}.

\printbibliography[heading=none]

Wikipedia \textbf{nu} este o resursă bibliografică acceptabilă.

\newpage


\section{Reguli de citare}

\textbf{Este foarte puternic recomandat} să nu se preia ca atare fragmente de text din alte lucrări, ci să fie reformulate ideile, în particular într-un mod sintetic și cu accent pe relevanța lor pentru lucrarea de disertație.

Dacă, totuși, se preia ca atare un fragment de text din altă lucrare, atunci este \textbf{obligatoriu} ca acesta să fie scris între \textbf{ghilimele} și să includă \textbf{referința bibliografică} – ex.: „Modelul sursă-filtru exprimă faptul că efectele compuse ale tractului vocal și radiației buzelor sunt reprezentate printr-un filtru digital cu parametri lent variabili în timp.” [7].\\

\textbf{Este foarte puternic recomandat} ca toate tabelele și figurile prezentate în lucrarea de disertație să fie \textbf{originale}.

Dacă, totuși, se preia ca atare un tabel sau o figură din altă lucrare, atunci este \textbf{obligatoriu} ca legenda acestuia să includă \textbf{referința bibliografică} \textbf{--} ex.: \textbf{Figura 4.5.} Arhitectura rețelei AlexNet [8].

Dacă se preia un tabel sau o figură din altă lucrare și se fac doar modificări minore, atunci este \textbf{obligatoriu} ca legenda acestuia să includă \textbf{referința bibliografică} după care au fost făcute modificările \textbf{--} ex.: \textbf{Figura 4.5.} Exemplu de rețea neurală convoluțională (după [8]).

Evident, în lucrarea de disertație, tabelul/figura va fi numerotat(ă) conform regulilor din capitolul 3 din prezentul ghid; nu este relevantă numerotarea tabelului/figurii din lucrarea originală.

\section{Asumarea proprietății intelectuale (\textit{copyright})}
Implicit, proprietatea intelectuală asupra lucrării (\textit{copyright}-ul) aparține UNSTPB.
În cazul în care absolventul dorește să pretindă drepturile de proprietate intelectuală
pentru propria persoană sau pentru o companie comercială (în cazul în care a existat o astfel
de colaborare), absolventul trebuie să contacteze \textbf{în timp util} coordonatorul științific.


